\documentclass[aspectratio=169]{beamer}
\input{header}
\coursecode{JKSSB}

\title{Firmware, BIOS, UEFI \& Drivers}
\subtitle{Lesson 46 --- The Boot Firmware Layer}
\author{JKSSB Computer Awareness}
\date{\today}

\begin{document}
\frame{\titlepage}

\begin{frame}{What Is Firmware?}
  \begin{itemize}
    \item \key{Firmware} is a special kind of software stored in
      \key{non-volatile} memory on the hardware itself.
    \item It gives a device its basic, low-level startup instructions.
    \item It runs \key{before} the operating system and starts the hardware.
    \item Firmware is programmed into a read-only memory (ROM) or a flash
      chip, not loaded from a file each time.
  \end{itemize}
  \begin{block}{The one idea}
    Firmware sits between hardware and higher-level software; without it the
    hardware does not know how to start.
  \end{block}
\end{frame}

\begin{frame}{Where Firmware Is Stored}
  \begin{itemize}
    \item Firmware is kept in \key{non-volatile} memory, so it survives when
      the power is off.
    \item Common locations: \key{ROM} (Read-Only Memory), PROM, EPROM, EEPROM,
      and \key{flash} memory.
    \item Because it is non-volatile, the startup instructions are ready the
      moment the machine is switched on.
    \item This is why a computer can begin working before any file is read
      from disk.
  \end{itemize}
  \begin{alertblock}{Trap alert}
    Firmware lives in ROM/flash, \key{not} in RAM. RAM is volatile and is
    empty at power-on.
  \end{alertblock}
\end{frame}

\begin{frame}{System Firmware: The BIOS}
  \begin{itemize}
    \item \key{BIOS} stands for \key{Basic Input/Output System}.
    \item It is the firmware stored on the motherboard of a PC.
    \item It runs \key{first} when the computer is switched on.
    \item It \key{initialises the hardware} and then \key{loads the operating
      system}.
  \end{itemize}
  \begin{block}{Exam fact}
    The firmware that initialises hardware and loads the OS is the BIOS
    (PYQ-28).
  \end{block}
\end{frame}

\begin{frame}{The Power-On Boot Sequence}
  \begin{itemize}
    \item The machine follows a fixed chain each time it starts:
      \begin{enumerate}
        \item Power on.
        \item Firmware (\key{BIOS/UEFI}) runs.
        \item Hardware initialisation and \key{POST}.
        \item Firmware finds and loads the \key{operating system}.
        \item The OS starts and loads \key{device drivers}.
        \item \key{Application} programs run.
      \end{enumerate}
  \end{itemize}
  \vspace{0.2cm}
  \muted{power-on $\rightarrow$ firmware $\rightarrow$ initialisation
  $\rightarrow$ OS $\rightarrow$ driver $\rightarrow$ application}
\end{frame}

\begin{frame}{POST: Power-On Self-Test}
  \begin{itemize}
    \item \key{POST} stands for \key{Power-On Self-Test}.
    \item It is one of the first things the BIOS does after power-on.
    \item It checks that key hardware (CPU, RAM, keyboard, disks) is present
      and working.
    \item If POST fails, the computer may beep or show an error instead of
      booting.
  \end{itemize}
  \begin{alertblock}{Trap alert}
    POST is a BIOS step and runs \key{before} the operating system, not
    after it.
  \end{alertblock}
\end{frame}

\begin{frame}{What the BIOS Does}
  \begin{itemize}
    \item Runs the \key{POST} to test the basic hardware.
    \item \key{Initialises} and configures the basic hardware.
    \item Reads settings such as the \key{boot order}.
    \item Finds the boot device and hands control to the \key{operating
      system}.
    \item Provides basic input/output services to the OS at an early stage.
  \end{itemize}
\end{frame}

\begin{frame}{BIOS Setup and CMOS}
  \begin{itemize}
    \item BIOS settings --- date/time, boot order, passwords --- are stored in
      a small battery-backed chip called \key{CMOS}.
    \item The BIOS Setup utility is opened with a key during startup,
      commonly \key{Del} or \key{F2}.
    \item Keeping these settings on a battery-backed chip is why the clock
      and boot order are remembered between sessions.
  \end{itemize}
  \vspace{0.2cm}
  \muted{The exact setup key varies by manufacturer.}
\end{frame}

\begin{frame}{UEFI: The Modern Firmware}
  \begin{itemize}
    \item \key{UEFI} stands for \key{Unified Extensible Firmware Interface}.
    \item It is the modern replacement for the legacy BIOS.
    \item It starts the machine \key{faster}, supports large disks (GPT), and
      offers a graphical setup utility.
    \item It adds \key{Secure Boot}, which helps block unauthorised software
      from loading at startup.
  \end{itemize}
  \begin{block}{For the exam}
    The firmware that boots the machine may be called BIOS or UEFI depending
    on the age of the system.
  \end{block}
\end{frame}

\begin{frame}{BIOS vs UEFI}
  \begin{center}
    \begin{tabular}{p{0.16\textwidth}p{0.34\textwidth}p{0.40\textwidth}}
      \toprule
      \textbf{Point} & \textbf{BIOS} & \textbf{UEFI} \\
      \midrule
      Full form & Basic Input/Output System & Unified Extensible Firmware Interface \\
      Age & Legacy (older) & Modern \\
      Boot speed & Slower & Faster \\
      Disk support & MBR (about 2 TB limit) & GPT (very large disks) \\
      Interface & Mostly text-based & Graphical \\
      Security & No Secure Boot & Secure Boot available \\
      \bottomrule
    \end{tabular}
  \end{center}
\end{frame}

\begin{frame}{Device Drivers}
  \begin{itemize}
    \item A \key{device driver} is \key{software} that lets the operating
      system talk to a specific hardware device.
    \item Each device --- printer, graphics card, network card --- needs its
      own driver.
    \item The driver translates general OS commands into device-specific
      instructions.
    \item Drivers are loaded \key{by the operating system}, not before it.
  \end{itemize}
  \begin{block}{Keep it straight}
    A driver is software, \key{not} hardware, and it works \key{after} the OS
    has started.
  \end{block}
\end{frame}

\begin{frame}{Drivers in the Boot Chain}
  \begin{itemize}
    \item Sequence: firmware boots, the \key{OS} starts, then the OS loads
      the \key{drivers}.
    \item \key{BIOS/UEFI} starts the machine; drivers run only once the OS is
      running.
    \item A driver is matched to the exact device and OS version.
  \end{itemize}
  \begin{alertblock}{Trap alert}
    The layer that first initialises hardware and loads the OS is
    \key{BIOS/UEFI}, never the device driver (TRAP-24).
  \end{alertblock}
\end{frame}

\begin{frame}{Firmware vs Driver vs Operating System}
  \begin{center}
    \begin{tabular}{p{0.16\textwidth}p{0.21\textwidth}p{0.20\textwidth}p{0.31\textwidth}}
      \toprule
      \textbf{Layer} & \textbf{Stored} & \textbf{Runs} & \textbf{Role} \\
      \midrule
      Firmware & On the hardware (ROM/flash) & First, at power-on & Starts the machine; loads the OS \\
      Device driver & On disk, with the OS & After the OS starts & Talks to one device \\
      Operating system & On disk & After firmware & Manages resources and the user interface \\
      \bottomrule
    \end{tabular}
  \end{center}
  \vspace{0.25cm}
  \muted{Keep the three roles separate --- a favourite exam confusion.}
\end{frame}

\begin{frame}{Firmware Updates (Flashing)}
  \begin{itemize}
    \item A \key{firmware update} replaces the firmware stored on a device.
    \item The process is commonly called \key{flashing}.
    \item Updates fix bugs, add features, and patch security flaws.
    \item Firmware updates are much \key{rarer} than ordinary software
      updates.
  \end{itemize}
  \begin{block}{Best practice}
    Always use the manufacturer's update for the exact model and keep a
    stable power supply during the update.
  \end{block}
\end{frame}

\begin{frame}{Risks of a Firmware Update}
  \begin{itemize}
    \item Firmware is low-level, so a failed update can leave a device
      unusable.
    \item Interrupting power during flashing can ``brick'' the device.
    \item Loading the \key{wrong firmware} for the model can damage it.
  \end{itemize}
  \begin{alertblock}{Warning}
    Never switch off or restart the machine in the middle of a firmware
    update.
  \end{alertblock}
\end{frame}

\begin{frame}{Worked Example}
  \begin{exampleblock}{A PC reports ``no boot device found''}
    Follow the boot chain. The firmware (BIOS/UEFI) runs, then the POST
    passes, then the firmware looks for a boot device. It cannot find the
    operating system, so it stops and shows the message. The failing layer is
    the \key{firmware/boot configuration}, not the application.
  \end{exampleblock}
  \vspace{0.2cm}
  \muted{Walk the chain in order: firmware $\rightarrow$ POST $\rightarrow$ OS
  $\rightarrow$ driver $\rightarrow$ application.}
\end{frame}

\begin{frame}{Trap Alert: Keep the Boot Roles Straight}
  \begin{itemize}
    \item Firmware that initialises hardware and loads the OS is
      \key{BIOS/UEFI}, not a driver or kernel (TRAP-24, PYQ-28).
    \item A \key{device driver} is software; the OS is system software
      (TRAP-10).
    \item Firmware is stored in ROM/flash (non-volatile); RAM is volatile.
    \item \key{POST} runs before the OS, not after.
    \item A driver is loaded \key{by the OS}, after the firmware.
  \end{itemize}
\end{frame}

\begin{frame}{Lesson 46: Recall}
  \begin{itemize}
    \item \key{Firmware} is software in non-volatile memory (ROM/flash) that
      starts and controls hardware.
    \item \key{BIOS} = Basic Input/Output System; it initialises hardware and
      loads the OS.
    \item \key{POST} = Power-On Self-Test, run by the BIOS before the OS.
    \item \key{UEFI} = Unified Extensible Firmware Interface; the modern
      firmware with fast boot, GPT, and Secure Boot.
    \item The startup chain is power-on $\rightarrow$ firmware
      $\rightarrow$ initialisation $\rightarrow$ OS $\rightarrow$ driver
      $\rightarrow$ application.
    \item A \key{device driver} is software loaded by the OS; firmware is not
      a driver.
    \item A firmware update is called \key{flashing}; interrupting it can
      damage the device.
  \end{itemize}
\end{frame}
\end{document}
